\section{Experimental Setup}
\label{sec:setup}

\subsection{Dataset and Input Protocol}
The primary study uses RML2016.10a, which contains eleven modulation classes
over SNR values from $-20$ to $18$ dB in 2-dB increments. Each record is
represented by a 128-sample IQ observation. Per-example RMS normalization is
applied before augmentation and feature derivation, using the fixed numerical
stabilizer recorded in the feature implementation. The class order and the
eleven-logit output interface are fixed for every condition.

\subsection{Data Splits and Audit}
Split A is the existing deterministic stratified 60/20/20 partition with a
recorded seed and index hash. Split B is an independently generated
stratified 60/20/20 partition with its own seed and index hash. Before model
training, the study records dataset and split hashes, class-by-SNR counts,
within-split train/validation/test overlap, duplicate checks, and membership
overlap between the two replications.

The source dictionary does not expose a source-group identifier for the
underlying waveform generation. The audit therefore reports exact index and
duplicate checks supported by the available metadata and does not infer
source-group independence. This distinction is retained when interpreting the
second split.

\subsection{Run Matrix and Epoch-Budget Pilot}
The primary BB1 grid contains four feature configurations, five width
multipliers, and three training seeds on Split A, for 60 planned records. The
confirmatory subset uses F2 and F7 at $0.50\times$, $1.00\times$, and
$2.00\times$ on both splits with
the five-seed set $\{42,17,123,2024,31\}$, for 60 records; the first three
Split-A seeds are reused only when their configuration and code hashes match.
This manuscript analyzes this confirmatory subset; the remaining main-grid
cells (F3, F5, $0.75\times$, and $1.50\times$) are not analyzed here.
Split A uses seed 42 and Split B uses seed 7 for split generation. The BB2
control was planned for F2/F7 at the low/high widths on Split A with three
seeds, for 12 directional-control records, but is not reported here.

Before the main budget is locked, a validation-only pilot evaluates F2/F7 at
$0.50\times$/$2.00\times$ under 200- and 300-epoch cosine schedules with seeds 42, 17, and
123. For each matched cell and seed, the decision compares validation
accuracy at the two final epochs. The 200-epoch budget is accepted only when
the paired interval is within the predeclared $\pm0.15$ percentage-point
equivalence band and no predefined low/high-SNR validation contrast exceeds
0.30 percentage points. Otherwise, the 300-epoch budget is used. Pilot records
remain separately labeled; a pilot run can be reused only when its split,
seed, code hash, and configuration hash match the selected main protocol.

\subsection{Training and Checkpointing}
All main-grid conditions use AdamW \cite{Loshchilov2019AdamW} with learning
rate $10^{-3}$, weight decay $10^{-2}$, batch size 512, cross-entropy with
0.05 label smoothing, and global gradient-norm clipping at 1.0. The learning
rate follows cosine annealing with $T_{\max}$ equal to the declared total
epoch budget, following the cosine-schedule formulation of SGDR
\cite{Loshchilov2017SGDR}. Early stopping is disabled, so the final epoch is
the primary checkpoint and the best-validation checkpoint is retained only as
a secondary diagnostic. The augmentation policy is applied to normalized raw I/Q before feature derivation.
Phase rotation is sampled with probability $0.75$; circular shifts with
probability $0.50$ use integer offsets from $-8$ to $8$ samples; additive
Gaussian noise is applied with probability $0.50$ and standard deviation
$0.025$; and amplitude scaling is applied with probability $0.50$ using a
Gaussian scale perturbation with standard deviation $0.08$, clipped to
$[0.75,1.25]$. The same operation order, automatic mixed-precision policy, and
random-seed handling are recorded for every run.

Every epoch checkpoint stores model, optimizer, scheduler, scaler, and random
number-generator states. A resumed run is linked to the interrupted record;
failed, interrupted, NaN, and manually stopped runs are retained in the event
log rather than overwritten. Validation predictions are written during the
training phase, while test predictions require the separate post-lock
inference command.

\subsection{Recorded Metrics and Outputs}
Each run records final and best validation accuracy, loss, per-SNR accuracy,
and region means. For the secondary SNR analysis, low, middle, and high
regions are defined as $\mathrm{SNR}<-8$ dB, $-8\leq\mathrm{SNR}\leq 2$ dB,
and $\mathrm{SNR}>2$ dB, respectively. The same records include per-class
accuracy and macro-F1, confusion
matrices, negative log-likelihood, expected calibration error, epoch curves,
training time, hardware utilization, and the complete configuration and
environment metadata. After test access is authorized, the same metrics are
computed for the final-epoch test predictions. Cost measurements are reported
separately as parameters, Core MACs, Core FLOPs, non-MAC FLOPs, preprocessing
FLOPs, and Total FLOPs.

The primary analysis is performed at run, seed, and split levels rather than
from a pooled sample-level accuracy alone. Generated tables and figures are
linked to the append-only records through their configuration, ledger, and
analysis hashes.
